%!TEX root = ../dokumentation.tex
%
% Nahezu alle Einstellungen koennen hier getaetigt werden
%
\RequirePackage[l2tabu, orthodox]{nag}  % weist in Commandozeile bzw. log auf veraltete LaTeX Syntax hin
 
\documentclass[%
    pdftex,
    oneside,            % Einseitiger Druck.
    12pt,               % Schriftgroesse
    parskip=half,       % Halbe Zeile Abstand zwischen Absätzen.
%   topmargin = 10pt,   % Abstand Seitenrand (Std:1in) zu Kopfzeile [laut log: unused]
    headheight = 12pt,  % Höhe der Kopfzeile
%   headsep = 30pt, % Abstand zwischen Kopfzeile und Text Body  [laut log: unused]
    headsepline,        % Linie nach Kopfzeile.
    footsepline,        % Linie vor Fusszeile.
    footheight = 16pt,  % Höhe der Fusszeile
    abstracton,     % Abstract Überschriften
    DIV=calc,       % Satzspiegel berechnen
    BCOR=8mm,       % Bindekorrektur links: 8mm
    headinclude=false,  % Kopfzeile nicht in den Satzspiegel einbeziehen
    footinclude=false,  % Fußzeile nicht in den Satzspiegel einbeziehen
    listof=totoc,       % Abbildungs-/ Tabellenverzeichnis im Inhaltsverzeichnis darstellen
    toc=bibliography,   % Literaturverzeichnis im Inhaltsverzeichnis darstellen
]{scrreprt} % Koma-Script report-Klasse, fuer laengere Bachelorarbeiten alternativ auch: scrbook
 
% Einstellungen laden
\usepackage{xstring}
\usepackage[utf8]{inputenc}
\usepackage[T1]{fontenc}
 
\newcommand{\einstellung}[1]{%
  \expandafter\newcommand\csname #1\endcsname{}
  \expandafter\newcommand\csname setze#1\endcsname[1]{\expandafter\renewcommand\csname#1\endcsname{##1}}
}
\newcommand{\langstr}[1]{\einstellung{lang#1}}
 
\input{ads/einstellungen_liste} % verfügbare Einstellungen
%%%%%%%%%%%%%%%%%%%%%%%%%%%%%%%%%%%%%%%%%%%%%%%%%%%%%%%%%%%%%%%%%%%%%%%%%%%%%%%
%                                   Einstellungen
%
% Hier können alle relevanten Einstellungen für diese Arbeit gesetzt werden.
% Dazu gehören Angaben u.a. über den Autor sowie Formatierungen.
%
%
%%%%%%%%%%%%%%%%%%%%%%%%%%%%%%%%%%%%%%%%%%%%%%%%%%%%%%%%%%%%%%%%%%%%%%%%%%%%%%%


%%%%%%%%%%%%%%%%%%%%%%%%%%%%%%%%%%%% Sprache %%%%%%%%%%%%%%%%%%%%%%%%%%%%%%%%%%%
%% Aktuell sind Deutsch und Englisch unterstützt.
%% Es werden nicht nur alle vom Dokument erzeugten Texte in
%% der entsprechenden Sprache angezeigt, sondern auch weitere
%% Aspekte angepasst, wie z.B. die Anführungszeichen und
%% Datumsformate.
\setzesprache{de} % oder en
%%%%%%%%%%%%%%%%%%%%%%%%%%%%%%%%%%%%%%%%%%%%%%%%%%%%%%%%%%%%%%%%%%%%%%%%%%%%%%%%

%%%%%%%%%%%%%%%%%%%%%%%%%%%%%%%%%%% Angaben  %%%%%%%%%%%%%%%%%%%%%%%%%%%%%%%%%%%
%% Die meisten der folgenden Daten werden auf dem
%% Deckblatt angezeigt, einige auch im weiteren Verlauf
%% des Dokuments.
\setzematrikelnr{1695387}
\setzekurs{T41000}
\setzetitel{Evaluation und Implementierung einer robusten Netzwerkarchitektur zur Echtzeit-Übertragung von Sensordaten zwischen einem ESP32-Client und einem Raspberry Pi-basierten IoT-Stack}
\setzedatumAbgabe{07.09.2026}
\setzefirma{Mercedes-Benz AG}
\setzefirmenort{Stuttgart}
\setzeabgabeort{Stuttgart}
\setzeabschluss{T1000}
\setzestudiengang{Informatik}
\setzedhbw{Stuttgart}
\setzebetreuer{Paul Karcher B. Eng}
\setzegutachter{}
\setzezeitraum{01.10.2015 - 07.09.2026} % Hier kann man auch z.B. 12 Wochen als Zeitraum angeben
\setzearbeit{Projektarbeit}
\setzeautor{Viktoria Backhaus}
%%%%%%%%%%%%%%%%%%%%%%%%%%%%%%%%%%%%%%%%%%%%%%%%%%%%%%%%%%%%%%%%%%%%%%%%%%%%%%%%

%%%%%%%%%%%%%%%%%%%%%%%%%%%% Literaturverzeichnis %%%%%%%%%%%%%%%%%%%%%%%%%%%%%%
%% Bei Fehlern während der Verarbeitung bitte in ads/header.tex bei der
%% Einbindung des Pakets biblatex (ungefähr ab Zeile 110,
%% einmal für jede Sprache), biber in bibtex ändern.
\newcommand{\ladeliteratur}{%
%\addbibresource{bibliographie.bib}
\addbibresource{quellen.bib}
}
%% Zitierstil
%% siehe: http://ctan.mirrorcatalogs.com/macros/latex/contrib/biblatex/doc/biblatex.pdf (3.3.1 Citation Styles)
%% mögliche Werte z.B numeric-comp, alphabetic, authoryear
\setzezitierstil{numeric-comp}
%%%%%%%%%%%%%%%%%%%%%%%%%%%%%%%%%%%%%%%%%%%%%%%%%%%%%%%%%%%%%%%%%%%%%%%%%%%%%%%%

%%%%%%%%%%%%%%%%%%%%%%%%%%%%%%%%% Layout %%%%%%%%%%%%%%%%%%%%%%%%%%%%%%%%%%%%%%%
%% Verschiedene Schriftarten
% laut nag Warnung: palatino obsolete, use mathpazo, helvet (option scaled=.95), courier instead
\setzeschriftart{lmodern} % palatino oder goudysans, lmodern, libertine

%% Paket um Textteile drehen zu können
%\usepackage{rotating}
%% Paket um Seite im Querformat anzuzeigen
%\usepackage{lscape}

%% Seitenränder
\setzeseitenrand{2.5cm}

%% Abstand vor Kapitelüberschriften zum oberen Seitenrand
\setzekapitelabstand{20pt}

%% Spaltenabstand
\setzespaltenabstand{10pt}
%%Zeilenabstand innerhalb einer Tabelle
\setzezeilenabstand{1.5}
%%%%%%%%%%%%%%%%%%%%%%%%%%%%%%%%%%%%%%%%%%%%%%%%%%%%%%%%%%%%%%%%%%%%%%%%%%%%%%%%

%%%%%%%%%%%%%%%%%%%%%%%%%%%%% Verschiedenes  %%%%%%%%%%%%%%%%%%%%%%%%%%%%%%%%%%%
%% Farben (Angabe in HTML-Notation mit großen Buchstaben)
\newcommand{\ladefarben}{%
	\definecolor{LinkColor}{HTML}{000000} % war ursprünglich 00007A
	\definecolor{ListingBackground}{HTML}{FCF7DE}
}
%% Mathematikpakete benutzen (Pakete aktivieren)
%\usepackage{amsmath}
%\usepackage{amssymb}

%% Programmiersprachen Highlighting (Listings)
\newcommand{\listingsettings}{%
	\lstset{%
		language=Java,			% Standardsprache des Quellcodes
		numbers=left,			% Zeilennummern links
		stepnumber=1,			% Jede Zeile nummerieren.
		numbersep=5pt,			% 5pt Abstand zum Quellcode
		numberstyle=\tiny,		% Zeichengrösse 'tiny' für die Nummern.
		breaklines=true,		% Zeilen umbrechen wenn notwendig.
		breakautoindent=true,	% Nach dem Zeilenumbruch Zeile einrücken.
		postbreak=\space,		% Bei Leerzeichen umbrechen.
		tabsize=2,				% Tabulatorgrösse 2
		basicstyle=\ttfamily\footnotesize, % Nichtproportionale Schrift, klein für den Quellcode
		showspaces=false,		% Leerzeichen nicht anzeigen.
		showstringspaces=false,	% Leerzeichen auch in Strings ('') nicht anzeigen.
		extendedchars=true,		% Alle Zeichen vom Latin1 Zeichensatz anzeigen.
		captionpos=b,			% sets the caption-position to bottom
		backgroundcolor=\color{ListingBackground}, % Hintergrundfarbe des Quellcodes setzen.
		xleftmargin=0pt,		% Rand links
		xrightmargin=0pt,		% Rand rechts
		frame=single,			% Rahmen an
		frameround=ffff,
		rulecolor=\color{darkgray},	% Rahmenfarbe
		fillcolor=\color{ListingBackground},
		keywordstyle=\color[rgb]{0.133,0.133,0.6},
		commentstyle=\color[rgb]{0.133,0.545,0.133},
		stringstyle=\color[rgb]{0.627,0.126,0.941}
	}
}
%%%%%%%%%%%%%%%%%%%%%%%%%%%%%%%%%%%%%%%%%%%%%%%%%%%%%%%%%%%%%%%%%%%%%%%%%%%%%%%%

%%%%%%%%%%%%%%%%%%%%%%%%%%%%%%%% Eigenes %%%%%%%%%%%%%%%%%%%%%%%%%%%%%%%%%%%%%%%
%% Hier können Ergänzungen zur Präambel vorgenommen werden (eigene Pakete, Einstellungen)
%% custom Quellenangabe
\usepackage{xifthen}
% first: indirect or direct
% second: lecture
% third: startpage
% fourth: endpage
\newcommand{\citeM}[4]
{(\ifthenelse{\equal{#1}{i}}{vgl.~}{}\cite{#2}\ifthenelse{\equal{#3}{} \AND \equal{#4}{}}{}{\ifthenelse{\equal{#1}{c}}{,~Z.}{,~S.}\ifthenelse{\equal{#3}{#4}}{#3}{#3-#4}})}

\usepackage{tikz}
\usetikzlibrary{shapes.geometric, arrows}

\usepackage{smartdiagram}


 % lese Einstellungen
 
\input{lang/strings} % verfügbare Strings
\input{lang/\sprache} % Übersetzung einlesen
 
% Einstellung der Sprache des Paketes Babel und der Verzeichnisüberschriften
\iflang{de}{\usepackage[english, ngerman]{babel}}
\iflang{en}{\usepackage[ngerman, english]{babel}}
 
 
%%%%%%% Package Includes %%%%%%%
 
\usepackage[margin=\seitenrand,foot=1cm]{geometry}  % Seitenränder und Abstände
\usepackage[activate]{microtype} %Zeilenumbruch und mehr
\usepackage[onehalfspacing]{setspace}
\usepackage{makeidx}
\usepackage[autostyle=true,german=quotes]{csquotes}
\usepackage{longtable}
\usepackage{enumitem}   % mehr Optionen bei Aufzählungen
\usepackage{graphicx}
\usepackage{booktabs}
\usepackage{pdfpages}   % zum Einbinden von PDFs
\usepackage{xcolor}     % für HTML-Notation
\usepackage{float}
\usepackage{array}
\usepackage{calc}       % zum Rechnen (Bildtabelle in Deckblatt)
\usepackage[right]{eurosym}
\usepackage{wrapfig}
\usepackage{pgffor} % für automatische Kapiteldateieinbindung
\usepackage[perpage, hang, multiple, stable]{footmisc} % Fussnoten
\usepackage[printonlyused]{acronym} % falls gewünscht kann die Option footnote eingefügt werden, dann wird die Erklärung nicht inline sondern in einer Fußnote dargestellt
\usepackage{listings}
 
% Notizen. Einsatz mit \todo{Notiz} oder \todo[inline]{Notiz}.
\usepackage[obeyFinal,backgroundcolor=yellow,linecolor=black]{todonotes}
% Alle Notizen ausblenden mit der Option "final" in \documentclass[...] oder durch das auskommentieren folgender Zeile
% \usepackage[disable]{todonotes}
 
%%%%%%%%%%%%%%%
%own packages
\usepackage{pgfplots}
\usepackage{tikz}
\usepackage{subcaption}
\usetikzlibrary{shapes,arrows,calc,fit,backgrounds,shapes.multipart,math,angles,quotes,patterns}
\pgfplotsset{compat=1.18}
\usepackage{amsmath, amssymb}
\usepackage{interval}
\usepackage{struktex}
 
 
% Kommentarumgebung. Einsatz mit \comment{}. Alle Kommentare ausblenden mit dem Auskommentieren der folgenden und dem aktivieren der nächsten Zeile.
\newcommand{\comment}[1]{\par {\bfseries \color{blue} #1 \par}} %Kommentar anzeigen
% \newcommand{\comment}[1]{} %Kommentar ausblenden
 
 
%%%%%% Configuration %%%%%
 
%% Anwenden der Einstellungen
 
\usepackage{\schriftart}
\ladefarben{}
 
% Titel, Autor und Datum
\title{\titel}
\author{\autor}
\date{\datum}
 
% PDF Einstellungen
\usepackage[%
    pdftitle={\titel},
    pdfauthor={\autor},
    pdfsubject={\arbeit},
    pdfcreator={pdflatex, LaTeX with KOMA-Script},
    pdfpagemode=UseOutlines,        % Beim Oeffnen Inhaltsverzeichnis anzeigen
    pdfdisplaydoctitle=true,        % Dokumenttitel statt Dateiname anzeigen.
    pdflang={\sprache},             % Sprache des Dokuments.
]{hyperref}
 
% (Farb-)einstellungen für die Links im PDF
\hypersetup{%
    colorlinks=true,        % Aktivieren von farbigen Links im Dokument
    linkcolor=LinkColor,    % Farbe festlegen
    citecolor=LinkColor,
    filecolor=LinkColor,
    menucolor=LinkColor,
    urlcolor=LinkColor,
    linktoc=all,            % der komplette text eines Eintrages im Table of Contents wird klickbar
    linktocpage=false,      % Nicht der Text sondern die Seitenzahlen in Verzeichnissen klickbar wenn true, hab ich aber deaktiviert
    bookmarksnumbered=true  % Überschriftsnummerierung im PDF Inhalt anzeigen.
}
% Workaround um Fehler in Hyperref, muss hier stehen bleiben
\usepackage{bookmark} %nur ein latex-Durchlauf für die Aktualisierung von Verzeichnissen nötig
 
% Schriftart in Captions etwas kleiner
\addtokomafont{caption}{\small}
 
% Literaturverweise (sowohl deutsch als auch englisch)
\iflang{de}{%
\usepackage[
    backend=bibtex,     % empfohlen. Falls biber Probleme macht: bibtex
    bibwarn=true,
    bibencoding=utf8,   % wenn .bib in utf8, sonst ascii
    sortlocale=de_DE,
    style=\zitierstil,
]{biblatex}
\DeclareFieldFormat{urldate}{\mkbibparens{zuletzt geprüft\addcolon\space#1}}
\DeclareBibliographyDriver{online}{%
  \usebibmacro{bibindex}%
  \usebibmacro{begentry}%
  \usebibmacro{author/editor}%
  \setunit{\labelnamepunct}\newblock
  \usebibmacro{title}%
  \newunit\newblock
  \usebibmacro{byauthor}%
  \newunit\newblock
  \usebibmacro{byeditor+others}%
  \newunit\newblock
  \usebibmacro{date}%
  \newunit\newblock
  \usebibmacro{url}%
  \newunit\newblock%
  \printfield{note}%
  \newunit\newblock
  \usebibmacro{addendum+pubstate}%
  \usebibmacro{pageref}%
  \usebibmacro{finentry}}
}
\iflang{en}{%
\usepackage[
    backend=biber,      % empfohlen. Falls biber Probleme macht: bibtex
    bibwarn=true,
    bibencoding=utf8,   % wenn .bib in utf8, sonst ascii
    sortlocale=en_US,
    style=\zitierstil,
]{biblatex}
}
 
\ladeliteratur{}
 
% Glossar
\usepackage[nonumberlist,toc]{glossaries}
 
%%%%%% Additional settings %%%%%%
 
% Hurenkinder und Schusterjungen verhindern
% http://projekte.dante.de/DanteFAQ/Silbentrennung
\clubpenalty = 10000 % schließt Schusterjungen aus (Seitenumbruch nach der ersten Zeile eines neuen Absatzes)
\widowpenalty = 10000 % schließt Hurenkinder aus (die letzte Zeile eines Absatzes steht auf einer neuen Seite)
\displaywidowpenalty=10000
 
% Bildpfad
\graphicspath{{images/}}
 
% Einige häufig verwendete Sprachen
\lstloadlanguages{PHP,Python,Java,C,C++,bash}
\listingsettings{}
% Umbennung des Listings
\renewcommand\lstlistingname{\langlistingname}
\renewcommand\lstlistlistingname{\langlistlistingname}
\def\lstlistingautorefname{\langlistingautorefname}
 
% Abstände in Tabellen
\setlength{\tabcolsep}{\spaltenabstand}
\renewcommand{\arraystretch}{\zeilenabstand}
 
% speziell für Tabellen
\usepackage{tabularx}
\usepackage{longtable}
% left fixed width:
\newcolumntype{L}[1]{>{\raggedright\arraybackslash}m{#1}}